\documentclass{scrartcl}
\usepackage{graphicx}  % required for inserting images

% for cyrillic and Bulgarian language
\usepackage[utf8]{inputenc}
\usepackage[T2A]{fontenc}
\usepackage[bulgarian]{babel}

% math packages
\usepackage{amsmath}
\usepackage{float}
\usepackage{amsfonts}
\usepackage{amssymb}
\usepackage{amsthm}
\usepackage{cancel}

\usepackage{ragged2e}  % adds FlushLeft

\usepackage{listings}
\usepackage{xcolor}

\definecolor{codegreen}{rgb}{0,0.6,0}
\definecolor{codegray}{rgb}{0.5,0.5,0.5}
\definecolor{codepurple}{rgb}{0.58,0,0.82}
\definecolor{backcolour}{rgb}{0.95,0.95,0.92}

\lstdefinestyle{mystyle}{
    language=Octave,
    backgroundcolor=\color{backcolour},   
    commentstyle=\color{codegreen},
    keywordstyle=\color{magenta},
    numberstyle=\tiny\color{codegray},
    stringstyle=\color{codepurple},
    basicstyle=\ttfamily\footnotesize,
    breakatwhitespace=false,         
    breaklines=true,
    captionpos=b,
    keepspaces=false,                 
    numbersep=5pt,
    showspaces=false,
    showstringspaces=false,
    showtabs=false,
    tabsize=2,
    columns=fullflexible,
}

\lstset{style=mystyle}

\usepackage[colorlinks=true, linkcolor=blue, urlcolor=cyan]{hyperref}

\newtheorem{theorem}{Теорема}
\newtheorem{definition}{Дефиниция}

\newtheorem{lemma}{Лема}
\newtheorem{corollary}{Следствие}
\newtheorem{remark}{Забележка}

\title{ДИС - Определен интеграл}
\author{Ивайло Андреев + gemini 3 Flash}
\date{16 май 2026}

\begin{document}
\maketitle

\tableofcontents

\newpage

\begin{FlushLeft}

\section{Разбиване на интервал и суми на Дарбу}

Въвеждането на понятието определен (Риманов) интеграл започва геометрично с идеята за пресмятане на лице на криволинеен трапец чрез апроксимация с правоъгълници. За целта дефинираме структура върху затворения интервал.

\begin{definition}[Разбиване на интервал]
Нека $[a, b]$ е краен и затворен интервал в $\mathbb{R}$ ($a < b$). Множеството от точки $\tau = \{x_0, x_1, x_2, \dots, x_n\}$, такива че:
\[ a = x_0 < x_1 < x_2 < \dots < x_{n-1} < x_n = b \]
се нарича \textbf{разбиване} на интервала $[a, b]$. 

Интервалите $[x_{i-1}, x_i]$ за $i = 1, 2, \dots, n$ се наричат \textbf{частични интервали} на разбиването $\tau$, а техните дължини означаваме с $\Delta x_i = x_i - x_{i-1}$.
\end{definition}

\begin{definition}[Диаметър на разбиване]
Числото $\delta(\tau)$, дефинирано като най-голямата от дължините на частичните интервали, се нарича \textbf{диаметър (стъпка)} на разбиването $\tau$:
\[ \delta(\tau) = \max_{1 \le i \le n} \Delta x_i \]
\end{definition}

Нека $f: [a, b] \to \mathbb{R}$ е \textbf{ограничена функция} в интервала $[a, b]$. Тъй като е ограничена в целия интервал, тя е ограничена и във всеки частичен интервал $[x_{i-1}, x_i]$. Следователно съществуват точните долна и горна граници:
\[ m_i = \inf_{x \in [x_{i-1}, x_i]} f(x), \quad M_i = \sup_{x \in [x_{i-1}, x_i]} f(x) \]

\begin{definition}[Суми на Дарбу]
За дадено разбиване $\tau$ на интервала $[a, b]$ и ограничена функция $f(x)$ дефинираме:
\begin{enumerate}
    \item \textbf{Малка сума на Дарбу} $s(f, \tau)$ (сбор от лицата на вписаните правоъгълници):
    \[ s(f, \tau) = \sum_{i=1}^{n} m_i \Delta x_i \]
    \item \textbf{Голяма сума на Дарбу} $S(f, \tau)$ (сбор от лицата на описаните правоъгълници):
    \[ S(f, \tau) = \sum_{i=1}^{n} M_i \Delta x_i \]
\end{enumerate}
Понеже $m_i \le M_i$ за всяко $i$, очевидно е изпълнено неравенството $s(f, \tau) \le S(f, \tau)$.
\end{definition}

\subsection{Поведение на сумите на Дарбу при добавяне на нови точки}

Казваме, че разбиването $\tau'$ е \textbf{наслагване (разширение)} на разбиването $\tau$ (означава се $\tau \subset \tau'$), ако $\tau'$ съдържа всички точки на $\tau$ и поне още една нова точка.

\begin{theorem}
При добавяне на нови точки в разбиването на интервала, големите суми на Дарбу не нарастват, а малките не намаляват. Тоест, ако $\tau \subset \tau'$, то:
\[ s(f, \tau) \le s(f, \tau') \quad \text{и} \quad S(f, \tau') \le S(f, \tau) \]
\end{theorem}

\begin{proof}
Достатъчно е да докажем твърдението за случая, когато към разбиването $\tau$ се добавя точно \textbf{една нова точка} $y$. Нека тази точка попадне в частичния интервал $[x_{k-1}, x_k]$, т.е. $x_{k-1} < y < x_k$.

Тогава новият интервал е разделен на два подобласта: $[x_{k-1}, y]$ и $[y, x_k]$. Означаваме точните граници на функцията в тези нови интервали:
\[ m_k' = \inf_{x \in [x_{k-1}, y]} f(x), \quad m_k'' = \inf_{x \in [y, x_k]} f(x) \]
\[ M_k' = \sup_{x \in [x_{k-1}, y]} f(x), \quad M_k'' = \sup_{x \in [y, x_k]} f(x) \]

Тъй като множествата $[x_{k-1}, y]$ и $[y, x_k]$ са подмножества на оригиналния интервал $[x_{k-1}, x_k]$, точният инфимум върху по-тясно множество е по-голям или равен на инфимума върху по-широкото множество, а супремумът е по-малък или равен. Следователно:
\[ m_k \le m_k' \quad \text{и} \quad m_k \le m_k'' \]
\[ M_k \ge M_k' \quad \text{и} \quad M_k \ge M_k'' \]

Разглеждаме разликата между новата малка сума $s(f, \tau')$ и старата $s(f, \tau)$. Тъй като всички останали частични интервали извън $k$-тия са идентични, техните членове в сумата се съкращават:
\[ s(f, \tau') - s(f, \tau) = m_k'(y - x_{k-1}) + m_k''(x_k - y) - m_k(x_k - x_{k-1}) \]
Заместваме идентичността $(x_k - x_{k-1}) = (y - x_{k-1}) + (x_k - y)$:
\[ s(f, \tau') - s(f, \tau) = m_k'(y - x_{k-1}) + m_k''(x_k - y) - m_k(y - x_{k-1}) - m_k(x_k - y) \]
\[ s(f, \tau') - s(f, \tau) = \underbrace{(m_k' - m_k)}_{\ge 0}\underbrace{(y - x_{k-1})}_{> 0} + \underbrace{(m_k'' - m_k)}_{\ge 0}\underbrace{(x_k - y)}_{> 0} \ge 0 \]
Следователно $s(f, \tau') - s(f, \tau) \ge 0 \Rightarrow s(f, \tau) \le s(f, \tau')$.

Аналогично разглеждаме разликата за големите суми:
\[ S(f, \tau) - S(f, \tau') = M_k(x_k - x_{k-1}) - \left[ M_k'(y - x_{k-1}) + M_k''(x_k - y) \right] \]
\[ S(f, \tau) - S(f, \tau') = \underbrace{(M_k - M_k')}_{\ge 0}\underbrace{(y - x_{k-1})}_{> 0} + \underbrace{(M_k - M_k'')}_{\ge 0}\underbrace{(x_k - y)}_{> 0} \ge 0 \]
Следователно $S(f, \tau) - S(f, \tau') \ge 0 \Rightarrow S(f, \tau') \le S(f, \tau)$.

Ако към разбиването се добавят $m$ на брой точки, твърдението се доказва индуктивно чрез прилагане на горната процедура стъпка по стъпка $m$ пъти.
\end{proof}

\begin{corollary}
Ако $\tau_1$ и $\tau_2$ са \textbf{произволни} две разбивания на интервала $[a, b]$, то малката сума на Дарбу за кое да е разбиване не надминава голямата сума за което и да е друго разбиване:
\[ s(f, \tau_1) \le S(f, \tau_2) \]
\end{corollary}
\begin{proof}
Дефинираме общото разширение $\tau^* = \tau_1 \cup \tau_2$. Тогава от Теорема 1 следва:
\[ s(f, \tau_1) \le s(f, \tau^*) \le S(f, \tau^*) \le S(f, \tau_2) \]
Което доказва твърдението.
\end{proof}

\newpage

\section{Дефиниция на Риманов интеграл чрез Дарбу}

От Следствие 1 виждаме, че множеството от всички малки суми $\{s(f, \tau)\}$ е ограничено отгоре от коя да е голяма сума, а множеството от големи суми $\{S(f, \tau)\}$ е ограничено отдолу. Съгласно аксиомата за непрекъснатост на реалните числа, съществуват следните точни граници:

\begin{definition}[Долен и горен интеграл на Дарбу]~
\begin{enumerate}
    \item Точната горна граница на всички малки суми се нарича \textbf{долен интеграл на Дарбу}:
    \[ \underline{\int_{a}^{b}} f(x)\,dx = \sup_{\tau} s(f, \tau) \]
    \item Точната долна граница на всички големи суми се нарича \textbf{горен интеграл на Дарбу}:
    \[ \overline{\int_{a}^{b}} f(x)\,dx = \inf_{\tau} S(f, \tau) \]
\end{enumerate}
Винаги е изпълнено: $\underline{\int_{a}^{b}} f(x)\,dx \le \overline{\int_{a}^{b}} f(x)\,dx$.
\end{definition}

\begin{definition}[Риманов интеграл]
Казваме, че ограничената функция $f(x)$ е \textbf{интегруема по Риман} в затворения интервал $[a, b]$, ако нейните долен и горен интеграли на Дарбу са равни. 

Общата им стойност се нарича \textbf{определен (Риманов) интеграл} на функцията $f(x)$ от $a$ до $b$ и се означава с:
\[ \int_{a}^{b} f(x)\,dx = \underline{\int_{a}^{b}} f(x)\,dx = \overline{\int_{a}^{b}} f(x)\,dx \]
\end{definition}

\newpage

\section{Критерий за интегруемост}

\begin{theorem}[Критерий за интегруемост на Дарбу]
Дадена ограничена функция $f(x)$ е интегруема по Риман в интервала $[a, b]$ тогава и само тогава, когато за всяко $\varepsilon > 0$ съществува разбиване $\tau$ такова, че:
\[ S(f, \tau) - s(f, \tau) < \varepsilon \]
\end{theorem}

\begin{proof}~
\begin{enumerate}
    \item \textbf{Необходимост ($\Rightarrow$):}
    Нека $f(x)$ е интегруема, което означава, че $\underline{\int_{a}^{b}} f(x)\,dx = \overline{\int_{a}^{b}} f(x)\,dx = I$.
    По дефиниция $I = \sup_\tau s(f, \tau)$, следователно за всяко $\varepsilon > 0$ съществува разбиване $\tau_1$, такова че:
    \[ s(f, \tau_1) > I - \frac{\varepsilon}{2} \]
    Аналогично, понеже $I = \inf_\tau S(f, \tau)$, съществува разбиване $\tau_2$, такова че:
    \[ S(f, \tau_2) < I + \frac{\varepsilon}{2} \]
    Разглеждаме общото разширение $\tau = \tau_1 \cup \tau_2$. Тогава от свойствата за монотонност (Теорема 1) имаме:
    \[ I - \frac{\varepsilon}{2} < s(f, \tau_1) \le s(f, \tau) \le S(f, \tau) \le S(f, \tau_2) < I + \frac{\varepsilon}{2} \]
    Оттук директно пресмятаме разликата:
    \[ S(f, \tau) - s(f, \tau) < \left(I + \frac{\varepsilon}{2}\right) - \left(I - \frac{\varepsilon}{2}\right) = \varepsilon \]
    
    \item \textbf{Достатъчност ($\Leftarrow$):}
    Нека за всяко $\varepsilon > 0$ е изпълнено $S(f, \tau) - s(f, \tau) < \varepsilon$. По дефиниция за долния и горния интеграл знаем, че:
    \[ s(f, \tau) \le \underline{\int_{a}^{b}} f(x)\,dx \le \overline{\int_{a}^{b}} f(x)\,dx \le S(f, \tau) \]
    Следователно разликата между горния и долния интеграл удовлетворява неравенството:
    \[ 0 \le \overline{\int_{a}^{b}} f(x)\,dx - \underline{\int_{a}^{b}} f(x)\,dx \le S(f, \tau) - s(f, \tau) < \varepsilon \]
    Тъй като неравенството $0 \le \overline{\int_{a}^{b}} f(x)\,dx - \underline{\int_{a}^{b}} f(x)\,dx < \varepsilon$ е изпълнено за абсолютно всяко произволно малко $\varepsilon > 0$, единствената математическа възможност е тази разлика да бъде нула:
    \[ \overline{\int_{a}^{b}} f(x)\,dx = \underline{\int_{a}^{b}} f(x)\,dx \]
    Което означава, че функцията е интегруема по Риман.
\end{enumerate}
Теоремата е доказана.
\end{proof}

\newpage

\section{Интегруемост на непрекъснатите функции}

За доказване на следващата фундаментална теорема се използва една основна теорема от топологията/анализа, която съгласно изискванията на конспекта се приема без доказателство:

\begin{theorem}[Теорема на Кантор за равномерната непрекъснатост]
Всяка функция, която е непрекъсната в краен и затворен интервал $[a, b]$, е равномерно непрекъсната в него. Това означава, че за всяко $\varepsilon_0 > 0$ съществува $\delta_0 > 0$ такова, че за всеки две точки $x', x'' \in [a, b]$, за които $|x' - x''| < \delta_0$, е изпълнено:
\[ |f(x') - f(x'')| < \varepsilon_0 \]
\end{theorem}

\begin{theorem}
Всяка непрекъсната функция $f(x)$ в крайния и затворен интервал $[a, b]$ е интегруема по Риман в него.
\end{theorem}

\begin{proof}
Нека изберем произволно число $\varepsilon > 0$. Понеже $f(x)$ е непрекъсната в $[a, b]$, съгласно Теоремата на Кантор тя е равномерно непрекъсната в този интервал.

Избираме стойност за $\varepsilon_0 = \frac{\varepsilon}{b - a} > 0$. Тогава от равномерната непрекъснатост следва, че съществува $\delta > 0$ такова, че за произволни $x', x'' \in [a, b]$ с $|x' - x''| < \delta$ е вярно:
\[ |f(x') - f(x'')| < \frac{\varepsilon}{b - a} \]

Конструираме произволно разбиване $\tau$ на интервала $[a, b]$, чийто диаметър е по-малък от намереното $\delta$, т.е. $\delta(\tau) < \delta$. Това означава, че за всеки частичен интервал $[x_{i-1}, x_i]$ неговата дължина е $\Delta x_i < \delta$.

Тъй като функцията $f(x)$ е непрекъсната във всеки затворен частичен интервал $[x_{i-1}, x_i]$, съгласно *Първата и Втората теореми на Вайерщрас* тя е ограничена и достига своя минимум $m_i$ и максимум $M_i$ в конкретни точки от този интервал. Нека:
\[ m_i = f(x_i'), \quad M_i = f(x_i'') \quad \text{където } x_i', x_i'' \in [x_{i-1}, x_i] \]

Тъй като и двете точки $x_i'$ и $x_i''$ принадлежат на един и същ частичен интервал, разстоянието между тях е строго по-малко от неговата дължина, а тя е по-малка от $\delta$:
\[ |x_i' - x_i''| \le \Delta x_i \le \delta(\tau) < \delta \]
Прилагаме свойството за равномерна непрекъснатост за тези две точки:
\[ M_i - m_i = f(x_i'') - f(x_i') < \frac{\varepsilon}{b - a} \]

Сега образуваме разликата между голямата и малката суми на Дарбу за разбиването $\tau$:
\[ S(f, \tau) - s(f, \tau) = \sum_{i=1}^{n} M_i \Delta x_i - \sum_{i=1}^{n} m_i \Delta x_i = \sum_{i=1}^{n} (M_i - m_i) \Delta x_i \]
Заместваме полученото по-горе неравенство $M_i - m_i < \frac{\varepsilon}{b - a}$:
\[ S(f, \tau) - s(f, \tau) < \sum_{i=1}^{n} \frac{\varepsilon}{b - a} \Delta x_i = \frac{\varepsilon}{b - a} \sum_{i=1}^{n} \Delta x_i \]
Тъй като сумата от дължините на всички частични интервали е точно равна на дължината на целия интервал, т.е. $\sum_{i=1}^{n} \Delta x_i = b - a$, получаваме:
\[ S(f, \tau) - s(f, \tau) < \frac{\varepsilon}{b - a} \cdot (b - a) = \varepsilon \]

Доказахме, че за всяко $\varepsilon > 0$ съществува разбиване $\tau$, за което $S(f, \tau) - s(f, \tau) < \varepsilon$. Съгласно Критерия на Дарбу (Теорема 2), функцията $f(x)$ е интегруема по Риман в $[a, b]$.
\end{proof}

\newpage

\section{Основни свойства на определения интеграл}

Тук изброяваме основните свойства на Римановия интеграл (съгласно изискването на конспекта, без доказателство):
\begin{enumerate}
    \item \textbf{Линейност:} Ако $f$ и $g$ са интегруеми в $[a, b]$, и $\alpha, \beta \in \mathbb{R}$, то $\alpha f + \beta g$ е интегруема и:
    \[ \int_{a}^{b} (\alpha f(x) + \beta g(x))\,dx = \alpha \int_{a}^{b} f(x)\,dx + \beta \int_{a}^{b} g(x)\,dx \]
    \item \textbf{Адитивност спрямо интервала:} Ако $f$ е интегруема в най-широкия интервал, съдържащ точките $a, b, c$, то:
    \[ \int_{a}^{b} f(x)\,dx = \int_{a}^{c} f(x)\,dx + \int_{c}^{b} f(x)\,dx \]
    \item \textbf{Интегриране на неравенства:} Ако $f(x) \le g(x)$ за всяко $x \in [a, b]$, то:
    \[ \int_{a}^{b} f(x)\,dx \le \int_{a}^{b} g(x)\,dx \]
\end{enumerate}

\subsection{Теорема за средните стойности в интегралното смятане}

\begin{theorem}[Теорема за средната стойност]
Ако функцията $f(x)$ е непрекъсната в затворения интервал $[a, b]$, то съществува точка $c \in [a, b]$, такава че:
\[ \int_{a}^{b} f(x)\,dx = f(c)(b - a) \]
\end{theorem}

\begin{proof}
Тъй като $f(x)$ е непрекъсната в $[a, b]$, по теоремата на Вайерщрас тя достига своя точен минимум $m$ и точен максимум $M$ в този интервал:
\[ m \le f(x) \le M \quad \forall x \in [a, b] \]
Прилагаме свойството за интегриране на неравенства върху горния израз в граници от $a$ до $b$:
\[ \int_{a}^{b} m\,dx \le \int_{a}^{b} f(x)\,dx \le \int_{a}^{b} M\,dx \]
Тъй като $m$ и $M$ са константи, техните интеграли са просто $m(b-a)$ и $M(b-a)$:
\[ m(b - a) \le \int_{a}^{b} f(x)\,dx \le M(b - a) \]
Тъй като $b > a$, дължината на интервала е положителна $(b - a) > 0$. Разделяме трите страни на неравенството на $(b - a)$:
\[ m \le \frac{1}{b - a}\int_{a}^{b} f(x)\,dx \le M \]

Означаваме средната стойност на интеграла с числото $\mu = \frac{1}{b - a}\int_{a}^{b} f(x)\,dx$. Имаме $m \le \mu \le M$.

Сега използваме друга фундаментална теорема за непрекъснати функции (**Теорема на Болцано-Вайерщрас за междинните стойности**): всяка непрекъсната функция в затворен интервал приема всички стойности между минимума си $m$ и максимума си $M$. 

Тъй като числото $\mu$ се намира в интервала $[m, M]$, то задължително съществува точка $c \in [a, b]$, за която:
\[ f(c) = \mu \quad \Rightarrow \quad f(c) = \frac{1}{b - a}\int_{a}^{b} f(x)\,dx \]
Умножавайки двете страни по $(b - a)$, получаваме окончателното твърдение:
\[ \int_{a}^{b} f(x)\,dx = f(c)(b - a) \]
Теоремата е доказана.
\end{proof}

\newpage

\section{Теорема на Нютон-Лайбниц}

Тази теорема установява дълбоката връзка между диференциалното смятане (намиране на производна) и интегралното смятане (намиране на лице). За целта въвеждаме понятието "интеграл като функция на горната граница".

\begin{theorem}[Теорема за производна на интеграл по горната граница]
Нека $f(x)$ е непрекъсната функция в $[a, b]$. Тогава функцията $F(x)$, дефинирана като:
\[ F(x) = \int_{a}^{x} f(t)\,dt \quad \text{за } x \in [a, b] \]
е диференцируема в $[a, b]$ и нейната производна удовлетворява равенството:
\[ F'(x) = f(x) \quad \forall x \in [a, b] \]
Това означава, че $F(x)$ е примитивна функция на $f(x)$.
\end{theorem}

\begin{proof}
Нека фиксираме произволна точка $x \in [a, b]$ и и дадем нарастване $h \neq 0$ такова, че $(x + h) \in [a, b]$. По дефиниция за производна, трябва да пресметнем границата на диференчното частно при $h \to 0$:
\[ F(x + h) - F(x) = \int_{a}^{x+h} f(t)\,dt - \int_{a}^{x} f(t)\,dt \]
Използваме свойството за адитивност на интеграла $\int_{a}^{x+h} = \int_{a}^{x} + \int_{x}^{x+h}$:
\[ F(x + h) - F(x) = \left( \int_{a}^{x} f(t)\,dt + \int_{x}^{x+h} f(t)\,dt \right) - \int_{a}^{x} f(t)\,dt = \int_{x}^{x+h} f(t)\,dt \]

Тъй като $f(t)$ е непрекъсната в частичния интервал с краища $x$ и $x+h$, прилагаме **Теоремата за средната стойност** (Теорема 4) за този интеграл. Следователно съществува точка $c_h$ между $x$ и $x+h$, такава че:
\[ \int_{x}^{x+h} f(t)\,dt = f(c_h) \cdot \left[ (x + h) - x \right] = f(c_h) \cdot h \]
Тогава диференчното частно става:
\[ \frac{F(x + h) - F(x)}{h} = \frac{f(c_h) \cdot h}{h} = f(c_h) \]

Сега извършваме граничен преход при $h \to 0$. Тъй като точката $c_h$ е притисната между $x$ и $x+h$ (т.е. $x \le c_h \le x+h$ за $h>0$), по *Теоремата за трите делена (полицаите)*, когато $h \to 0$, точката $c_h \to x$.

Понеже по условие функцията $f$ е непрекъсната, можем да вкараме границата в аргумента:
\[ F'(x) = \lim_{h \to 0} \frac{F(x + h) - F(x)}{h} = \lim_{h \to 0} f(c_h) = f\left(\lim_{h \to 0} c_h\right) = f(x) \]
Доказахме, че $F'(x) = f(x)$, което означава, че интегралът като функция на горната граница е примитивна функция на подинтегралната функция.
\end{proof}

\begin{theorem}[Формула на Нютон-Лайбниц]
Нека $f(x)$ е непрекъсната функция в $[a, b]$ и $\Phi(x)$ е \textbf{произволна} нейна примитивна функция в този интервал (т.е. $\Phi'(x) = f(x)$). Тогава:
\[ \int_{a}^{b} f(x)\,dx = \Phi(b) - \Phi(a) \]
\end{theorem}

\begin{proof}
От Теорема 5 знаем, че функцията $F(x) = \int_{a}^{x} f(t)\,dt$ е конкретна примитивна функция на $f(x)$ в $[a, b]$. 

От основна теорема в диференциалното смятане (следствие от Лагранж) знаем, че кои да е две примитивни функции на една и съща подинтегрална функция се различават най-много с константа $C \in \mathbb{R}$. Тъй като по условие $\Phi(x)$ е някаква произволна примитивна функция, то:
\[ F(x) = \Phi(x) + C \quad \Rightarrow \quad \int_{a}^{x} f(t)\,dt = \Phi(x) + C \]

За да определим неизвестната константа $C$, заместваме в горното равенство стойността $x = a$:
\[ \int_{a}^{a} f(t)\,dt = \Phi(a) + C \]
Тъй като интеграл с еднакви граници е равен на нула ($\int_{a}^{a} = 0$), получаваме:
\[ 0 = \Phi(a) + C \quad \Rightarrow \quad C = -\Phi(a) \]

Заместваме намерената стойност за константата $C$ обратно в общото уравнение:
\[ \int_{a}^{x} f(t)\,dt = \Phi(x) - \Phi(a) \]
Сега, за да намерим крайния определен интеграл в граници от $a$ до $b$, просто заместваме променливата $x$ с горната граница $b$:
\[ \int_{a}^{b} f(t)\,dt = \Phi(b) - \Phi(a) \]
Тъй като променливата на интегриране е "променлива-фантом" (dummy variable), можем да сменим буквата обратно на $x$:
\[ \int_{a}^{b} f(x)\,dx = \Phi(b) - \Phi(a) \]
С това фундаменталната формула на Нютон-Лайбниц е напълно доказана.
\end{proof}

\end{FlushLeft}

\end{document}